\documentclass[11pt]{article}
\usepackage{amsmath,amssymb,amsthm}
\usepackage{geometry}
\usepackage{hyperref}
\usepackage{bm}
\usepackage{listings}
\usepackage{xcolor}
\geometry{margin=1in}

\lstset{
  basicstyle=\ttfamily\small,
  breaklines=true,
  frame=single,
  backgroundcolor=\color{gray!10},
  showstringspaces=false,
}

\newcommand{\dive}{\nabla\!\cdot\!}
\newcommand{\grad}{\nabla}
\newcommand{\nn}{\mathbf{n}}
\newcommand{\uu}{\mathbf{u}}
\newcommand{\II}{\mathbf{I}}
\newcommand{\SK}{\boldsymbol{\Sigma}_K}
\newcommand{\dxn}{\partial_n}
\newcommand{\Wb}{W_{\!\mathrm{wb}}}

\title{Well-Balanced Korteweg Capillary Stress for Diffuse Two-Phase
Flow: \\ Two Equivalent Forms and Their Discrete Cancellation Requirements}
\author{FLAMES / Alamo working notes}
\date{2026-05-15}

\begin{document}
\maketitle

\section{Goal}

The FLAMES Hydro2 inviscid path adds a Korteweg flux-form contribution
$-\SK\!\cdot\!\nn$ to the per-face HLLC momentum flux and
$-\SK\!\cdot\!\uu$ to the energy flux
(\texttt{Solver/Local/FluidRiemann/HLLC.H:263--284}). The intent is that on
the IC tanh, which is set up to be at Cahn--Hilliard (CH) equilibrium
($\lambda=2\kappa/\epsilon^2$ for the input convention $\eta=\tfrac12[1
+\tanh(x/\epsilon)]$), the Korteweg flux contribution to a planar interface
should be \emph{well-balanced} --- producing no net force, no net energy
work, and no observable interface signature in either pressure or velocity.

This note works through two algebraically distinct stress tensors that
appear in the literature, derives the well-balanced trace correction
required by each, and identifies the discrete cancellation property the
implementation must preserve so that no spurious pressure or velocity
features appear at the diffuse interface. We then map the analysis back
to the actual code in HLLC.H and Hydro2.cpp.

\paragraph{Convention.} Throughout, $W(\eta)=\eta^2(1-\eta)^2$ is the
standard symmetric double-well; $\lambda$ is the well height
(\texttt{ch\_W\_scale}); $\kappa$ is the gradient-energy coefficient
(\texttt{ch\_kappa\_eff}); $\mu_{\mathrm{chem}}\equiv-\kappa\grad^2\eta+\lambda W'(\eta)$
is the CH chemical potential. The Cahn--Hilliard equilibrium is
$\mu_{\mathrm{chem}}\equiv\mathrm{const}$. With $\mu_{\mathrm{chem}}=0$,
the planar tanh profile in the input convention satisfies
\begin{equation}
\kappa\,\grad^2\eta \;=\; \lambda\,W'(\eta), \qquad
\eta(x) = \tfrac12\!\left[1+\tanh\!\big(x/\epsilon\big)\right],
\qquad \lambda = \frac{2\kappa}{\epsilon^2}.
\label{eq:eq_relation}
\end{equation}

\section{Two candidate Korteweg stresses}

\subsection{Form~A --- classical binary-fluid (Anderson 1998 Eq.~19, no
\texorpdfstring{$\rho\grad^2\rho$}{rho grad squared rho} term)}

The form used in our theory writeup
(\texttt{theory/SQR\_source\_term\_derivation.tex}, Eq.~Sigmak):
\begin{equation}
\SK^{(A)} \;=\; \kappa\,\Big[\tfrac{1}{2}|\grad\eta|^2\,\II \;-\; \grad\eta\!\otimes\!\grad\eta\Big].
\label{eq:SigmaA}
\end{equation}
The components in a face-aligned $(n,t)$ frame are
\begin{equation}
\SK^{(A)}_{nn} = \tfrac{1}{2}\kappa|\grad\eta|^2 - \kappa(\dxn\eta)^2,
\qquad
\SK^{(A)}_{nt} = -\kappa(\dxn\eta)(\partial_t\eta).
\label{eq:SigmaAcomp}
\end{equation}

\subsection{Form~B --- ``thermo-consistent'' / Anderson 1998 Eq.~16
\texorpdfstring{$+$}{plus} CH-pressure}

The form actually implemented in \texttt{HLLC.H:276}, written without the
trace correction:
\begin{equation}
\SK^{(B)} \;=\; \kappa\,\Big[\tfrac{1}{2}|\grad\eta|^2\,\II \;-\; \grad\eta\!\otimes\!\grad\eta\Big] \;+\; \kappa\,\eta\,\grad^2\eta\,\II.
\label{eq:SigmaB}
\end{equation}
The extra diagonal term $+\kappa\eta\grad^2\eta\,\II$ does not appear in
the binary-fluid Korteweg stress used in our \S5 pillbox argument. Its
provenance is Cahn--Hilliard / Anderson Eq.~16: a $\rho\grad^2\rho$
analogue rewritten in $\eta$-variables when the gradient-energy is
written as $\tfrac{1}{2}\kappa|\grad\eta|^2$ inside the free energy
functional.

Components:
\begin{equation}
\SK^{(B)}_{nn} = \tfrac{1}{2}\kappa|\grad\eta|^2 - \kappa(\dxn\eta)^2 + \kappa\,\eta\,\grad^2\eta,
\qquad
\SK^{(B)}_{nt} = -\kappa(\dxn\eta)(\partial_t\eta).
\label{eq:SigmaBcomp}
\end{equation}
The off-diagonal $\SK_{nt}$ is identical between the two forms.

\section{Body force from the divergence}

We need two vector calculus identities:
\begin{align}
\dive(\grad\eta\!\otimes\!\grad\eta)
   &= (\grad^2\eta)\grad\eta + \tfrac{1}{2}\grad|\grad\eta|^2,
\label{eq:id1}\\
\dive\!\big(\phi\,\II\big) &= \grad\phi.
\label{eq:id2}
\end{align}

\subsection{Form~A}

Applying \eqref{eq:id1}--\eqref{eq:id2} to \eqref{eq:SigmaA}:
\begin{equation}
\dive\SK^{(A)}
   = \tfrac{1}{2}\kappa\grad|\grad\eta|^2 - \kappa(\grad^2\eta)\grad\eta - \tfrac{1}{2}\kappa\grad|\grad\eta|^2
   = -\,\kappa(\grad^2\eta)\grad\eta.
\label{eq:divA}
\end{equation}

This is the potential form $\fk = -\mu_K\grad\eta$ with $\mu_K=\kappa\grad^2\eta$
(the gradient-energy piece of $\mu_{\mathrm{chem}}$, opposite sign).

\subsection{Form~B}

Form~B differs from Form~A by $\kappa\eta\grad^2\eta\,\II$, whose divergence is
\begin{equation}
\dive(\kappa\eta\grad^2\eta\,\II) = \kappa\,\grad(\eta\grad^2\eta)
   = \kappa(\grad\eta)(\grad^2\eta) + \kappa\,\eta\,\grad(\grad^2\eta).
\label{eq:divextra}
\end{equation}
Adding to \eqref{eq:divA}: the $-\kappa(\grad^2\eta)\grad\eta$ from Form~A
\emph{cancels exactly} with the $+\kappa(\grad\eta)(\grad^2\eta)$ from
\eqref{eq:divextra} (they are the same vector with opposite sign), leaving
\begin{equation}
\boxed{\,\dive\SK^{(B)} \;=\; \kappa\,\eta\,\grad(\grad^2\eta).\,}
\label{eq:divB}
\end{equation}
This is \emph{not} the binary-fluid body force from \S5 of the theory
writeup; it is the gradient-energy piece of $-\eta\grad\mu_{\mathrm{chem}}$:
\begin{equation}
\dive\SK^{(B)} = -\,\eta\,\grad\!\big[-\kappa\grad^2\eta\big]
              = -\,\eta\,\grad\mu_K^{\!(B)},\qquad \mu_K^{\!(B)}\equiv-\kappa\grad^2\eta.
\label{eq:divBalt}
\end{equation}

\paragraph{Why this distinction matters.} At a planar interface in CH
equilibrium ($\grad^2\eta\ne 0$ pointwise inside the band but
$\mu_{\mathrm{chem}}=\mathrm{const}$), Form~A's body force is
$-\kappa(\grad^2\eta)\grad\eta = -\lambda W'(\eta)\grad\eta = -\grad[\lambda W(\eta)]$ ---
a gradient, but \emph{not zero}. Form~B's body force is
$\kappa\eta\grad(\grad^2\eta)$ which equals $\eta\,\lambda W''(\eta)\grad\eta$
at equilibrium (using \eqref{eq:eq_relation}). Neither form is
``well-balanced'' on its own. A trace correction must be added to each
so the residual body force vanishes at planar CH equilibrium.

\section{Bulk-thermo trace correction}

We choose a scalar $\Wb(\eta)$ and define
\begin{equation}
\SK^{\!\mathrm{wb}} \;=\; \SK \;-\; \Wb(\eta)\,\II,
\label{eq:SKwb}
\end{equation}
with the sign convention that the diagonal of $\SK^{\!\mathrm{wb}}$ is
the original diagonal \emph{minus} $\Wb$. The divergence picks up
$-\grad\Wb = -\Wb'(\eta)\grad\eta$. The well-balanced requirement is
\begin{equation}
\dive\SK^{\!\mathrm{wb}}\big|_{\mathrm{CH\;eq}} = 0.
\label{eq:wb}
\end{equation}

\subsection{Form~A: \texorpdfstring{$\Wb^{(A)}(\eta) = -\lambda W(\eta)$}{Wb(A) = -lambda W}}

Using \eqref{eq:divA} and \eqref{eq:eq_relation}:
\begin{equation}
\dive\SK^{(A,\mathrm{wb})}
   = -\kappa(\grad^2\eta)\grad\eta - \Wb'(\eta)\grad\eta
   = -[\kappa\grad^2\eta + \Wb'(\eta)]\grad\eta.
\end{equation}
At CH equilibrium $\kappa\grad^2\eta = \lambda W'(\eta)$, so the bracket
vanishes iff $\Wb'(\eta) = -\lambda W'(\eta)$, i.e.\
$\Wb^{(A)}(\eta) = -\lambda W(\eta) + \mathrm{const}$. Dropping the constant
(absorbed into the bulk pressure):
\begin{equation}
\boxed{\,\Wb^{(A)}(\eta) = -\lambda W(\eta),\qquad
\SK^{(A,\mathrm{wb})} = \kappa\big[\tfrac12|\grad\eta|^2\II - \grad\eta\!\otimes\!\grad\eta\big] + \lambda W(\eta)\,\II.\,}
\label{eq:Awb}
\end{equation}

\subsection{Form~B: \texorpdfstring{$\Wb^{(B)}(\eta) = \lambda[\eta W'(\eta)-W(\eta)]$}{Wb(B) = lambda(eta W' - W)}}

Using \eqref{eq:divB}:
\begin{equation}
\dive\SK^{(B,\mathrm{wb})} = \kappa\eta\grad(\grad^2\eta) - \Wb'(\eta)\grad\eta.
\end{equation}
At CH equilibrium, $\kappa\grad^2\eta = \lambda W'(\eta)$ so
$\grad(\kappa\grad^2\eta) = \lambda W''(\eta)\grad\eta$, hence
$\kappa\eta\grad(\grad^2\eta) = \lambda\eta W''(\eta)\grad\eta$.
The well-balanced requirement becomes
\begin{equation}
\Wb'(\eta) = \lambda\eta W''(\eta).
\end{equation}
Integrating: $\Wb(\eta) = \lambda\eta W'(\eta) - \lambda W(\eta) + \mathrm{const}$
(integration by parts), so
\begin{equation}
\boxed{\,\Wb^{(B)}(\eta) = \lambda\big[\eta W'(\eta) - W(\eta)\big],\qquad
\SK^{(B,\mathrm{wb})} = \SK^{(B)} - \Wb^{(B)}(\eta)\,\II.\,}
\label{eq:Bwb}
\end{equation}
For the double-well $W(\eta)=\eta^2(1-\eta)^2$:
\begin{equation}
\Wb^{(B)}(\eta) = \lambda\,\eta^2(3\eta-1)(\eta-1).
\label{eq:Wb_explicit}
\end{equation}
\textbf{This is exactly the form coded as} \texttt{q\_eta} \textbf{at} \texttt{Hydro2.cpp:1780--1782}\textbf{:}
\begin{lstlisting}
auto Q_global = [=] AMREX_GPU_DEVICE (Real e) -> Real {
    return lambda_W_ * e * e * (Real(3.0)*e - Real(1.0)) * (e - Real(1.0));
};
\end{lstlisting}
and \texttt{HLLC.H:276} encodes the full WB stress as
\begin{lstlisting}
Set::Scalar S_nn = -k*dn*dn + k*etaf*lapf + 0.5*k*ge2 - q;
\end{lstlisting}
i.e.\ $\SK^{(B,\mathrm{wb})}_{nn} = -\kappa(\dxn\eta)^2 + \kappa\eta\grad^2\eta + \tfrac12\kappa|\grad\eta|^2 - \lambda(\eta W' - W)$.

\subsection{The two well-balanced forms differ by an irrotational piece}

By construction both satisfy $\dive\SK^{(\cdot,\mathrm{wb})}=0$ at CH
equilibrium. Their difference is purely diagonal:
\begin{align}
\SK^{(B,\mathrm{wb})} - \SK^{(A,\mathrm{wb})}
   &= \kappa\eta\grad^2\eta\,\II - \big[\Wb^{(B)} - \Wb^{(A)}\big]\II \nonumber\\
   &= \kappa\eta\grad^2\eta\,\II - \big[\lambda(\eta W' - W) + \lambda W\big]\II \nonumber\\
   &= \kappa\eta\grad^2\eta\,\II - \lambda\eta W'(\eta)\,\II \nonumber\\
   &= \eta\big[\kappa\grad^2\eta - \lambda W'(\eta)\big]\II \nonumber\\
   &= -\,\eta\,\mu_{\mathrm{chem}}\,\II.
\label{eq:diff}
\end{align}
At CH equilibrium $\mu_{\mathrm{chem}}\equiv\mathrm{const}$, so the
difference is a constant scalar times $\II$ at equilibrium (absorbed into
the pressure; if $\mu_{\mathrm{chem}}=0$, the two forms agree pointwise on
the band). Off equilibrium they differ by a real pressure-like term
$-\eta\mu_{\mathrm{chem}}\II$.

\section{Pointwise vanishing on the planar tanh (1D equipartition)}

The well-balanced property \eqref{eq:wb} says only that the
\emph{divergence} of the stress vanishes at equilibrium. For 1D
\emph{planar} interfaces we can prove a stronger result: the normal
component of the stress vanishes \emph{pointwise} on the band.

\subsection{The CH equipartition identity}

For the 1D planar profile satisfying \eqref{eq:eq_relation} with
$\mu_{\mathrm{chem}}=0$, multiplying $\kappa\eta''=\lambda W'(\eta)$ by
$\eta'$ and integrating once (using $\eta'\to 0$ at $x\to\pm\infty$):
\begin{equation}
\boxed{\,\tfrac{1}{2}\kappa(\eta')^2 \;=\; \lambda\,W(\eta) \quad\text{(equipartition).}\,}
\label{eq:equipart}
\end{equation}
This is the standard Cahn--Hilliard result: at the planar minimizer the
gradient energy density equals the bulk double-well energy density pointwise.

\subsection{Form~A: \texorpdfstring{$\SK^{(A,\mathrm{wb})}_{nn}=0$ pointwise}{normal stress vanishes pointwise}}

In 1D, $|\grad\eta|^2 = (\eta')^2 = (\dxn\eta)^2$. From \eqref{eq:Awb}:
\begin{align}
\SK^{(A,\mathrm{wb})}_{nn}
   &= -\kappa(\eta')^2 + \tfrac12\kappa(\eta')^2 + \lambda W(\eta) \nonumber\\
   &= -\tfrac12\kappa(\eta')^2 + \lambda W(\eta) \nonumber\\
   &\stackrel{\eqref{eq:equipart}}{=} -\lambda W(\eta) + \lambda W(\eta) \;=\; 0.
\label{eq:An0}
\end{align}

\subsection{Form~B: \texorpdfstring{$\SK^{(B,\mathrm{wb})}_{nn}=0$ pointwise}{normal stress vanishes pointwise}}

Substituting $\kappa\eta''=\lambda W'(\eta)$ from \eqref{eq:eq_relation}:
\begin{align}
\SK^{(B,\mathrm{wb})}_{nn}
   &= -\kappa(\eta')^2 + \kappa\eta\eta'' + \tfrac12\kappa(\eta')^2 - \lambda(\eta W' - W) \nonumber\\
   &= -\tfrac12\kappa(\eta')^2 + \lambda\eta W'(\eta) - \lambda\eta W'(\eta) + \lambda W(\eta) \nonumber\\
   &= -\tfrac12\kappa(\eta')^2 + \lambda W(\eta) \nonumber\\
   &\stackrel{\eqref{eq:equipart}}{=} 0.
\label{eq:Bn0}
\end{align}

\paragraph{Consequence.} On the IC tanh in 1D, every face-evaluated
$\SK^{(\cdot,\mathrm{wb})}_{nn}$ is \emph{identically zero in the
continuum}. The Korteweg momentum-flux modifier $-\SK\!\cdot\!\nn$ at every
face is zero. The Korteweg energy-flux modifier $-(\SK\!\cdot\!\uu)\!\cdot\!\nn$
is zero regardless of $\uu$, since it is $-\SK_{nn}\,u_n - \SK_{nt}\,u_t = 0
-0$ (in 1D $\SK_{nt}=0$ trivially).

\subsection{What planar 1D equipartition tells us about the dip}

If the IC is at CH equilibrium and the Korteweg flux modifier is zero
pointwise in continuum, the implementation must produce a nonzero
contribution only via \emph{one of the following sources of discrete
inconsistency}:

\begin{enumerate}
\item[(D1)] \textbf{Discrete equipartition violation.} The face stencils
for $|\grad\eta|^2$, $\eta\grad^2\eta$, and $\Wb^{(B)}(\eta)$ don't
combine to give discrete zero on the IC tanh. They will not, in general,
because the 9-point isotropic Laplacian, the central first-difference,
and the face-averaged $\eta$ have different leading truncation errors.
\item[(D2)] \textbf{IC tanh is not exactly at the discrete CH
minimizer.} The continuum relation $\lambda=2\kappa/\epsilon^2$ is
exact only for the continuum tanh; the discrete minimizer of the
Hydro2 CH operator differs at $O(dx^2/\epsilon^2)$.
\item[(D3)] \textbf{Coupling to the Riemann state through $S_*$.} The
energy flux uses the HLLC contact velocity $S_*$. If
$\SK_{nn}\ne 0$ discretely on \emph{some} faces, the modifier
$-\SK_{nn}\,S_*$ varies with the local hydro state, and a passing wave
modulates the residual. This converts (D1)--(D2) into a
state-dependent energy artifact.
\end{enumerate}

(D3) is the path by which a small discrete residual in $\SK_{nn}$ ---
even at $O(\epsilon^2)$ truncation --- can be \emph{amplified} into an
observable pressure feature when a wave is present, because the
amplification factor is the wave's local velocity rather than the
asymptotically small interface velocity. This is consistent with the
empirical observation in
\texttt{input\_acoustic\_1d}: the dip appears only after the wave arrives
at the interface, not at $t=0$.

\section{Energy flux and the interstitial-working concern}

The Korteweg energy contribution at a face is
\begin{equation}
F_E^{(K)} = -\SK\!\cdot\!\uu\big|_{\mathrm{face}}\!\cdot\!\nn = -\SK_{nn}\,u_n - \SK_{nt}\,u_t.
\end{equation}
Its cell divergence in continuum is
\begin{equation}
\dive(\SK\!\cdot\!\uu) = (\dive\SK)\!\cdot\!\uu + \SK\!:\!\grad\uu.
\label{eq:edivflux}
\end{equation}
At CH equilibrium $(\dive\SK^{(\cdot,\mathrm{wb})}) = 0$, so the only
energy contribution from a perfectly well-balanced stress is the
\emph{Reynolds work} term $\SK\!:\!\grad\uu$.

In 1D planar, equipartition gives $\SK^{(\cdot,\mathrm{wb})}_{nn}=0$
pointwise (\S5.2--5.3), so $\SK^{(\cdot,\mathrm{wb})}\!:\!\grad\uu = 0$ as well.
\textbf{The Reynolds-work term is identically zero in continuum on the
planar CH-equilibrium band, even with a passing wave.} The dip cannot
come from the continuum Reynolds-work term --- only from its
\emph{discrete} residual.

\paragraph{Anderson interstitial working.} Anderson et al.\ (1998,
Eq.~16b) note that off-equilibrium 2nd-law consistency requires an
additional energy flux $\mathbf q_E^{\mathrm{int}} = -K_E (D\eta/Dt)\grad\eta$.
With \texttt{ch\_mobility\_nom}$=0$ (frozen $\eta$) and zero advection,
$D\eta/Dt = 0$ and the interstitial-working flux vanishes. This term is
\emph{not} the source of the dip in \texttt{input\_acoustic\_1d}.

\section{Discrete cancellation: what the implementation must preserve}

Translating the continuum analysis into a discrete fidelity criterion:

\begin{enumerate}
\item[(R1)] \textbf{Pointwise $\SK_{nn}$-zero on the IC tanh.} On the IC
profile evaluated at every face, the face stencils for
$(\dxn\eta)^2$, $|\grad\eta|^2$, $\eta\grad^2\eta$, and $\Wb(\eta)$
must combine to a value of $\SK_{nn}$ that is small enough that
$|\SK_{nn}^{\mathrm{disc}}|\!\cdot\!|u_{\mathrm{wave}}|$ stays below the
desired pressure tolerance once the wave has amplitude $u_{\mathrm{wave}}$.
\item[(R2)] \textbf{Telescoping cancellation between adjacent faces.}
Even if (R1) is only approximately satisfied, the discrete divergence
$(F_{n+}-F_{n-})/\Delta x$ should cancel to at least one order higher
than the per-face residual. This requires the same face stencil to be
used on both the $+$ and $-$ faces of a cell with no asymmetric
boundary treatment.
\item[(R3)] \textbf{State-independent face stress.} The Korteweg flux
modifier is computed from $\eta$ alone (lines 265--277 of
\texttt{HLLC.H}). This is correct --- but it means the face stress is
not symmetrized with the local hydro state, so any per-face residual
multiplies the local wave velocity through the energy update.
\end{enumerate}

\subsection{Face stencil in Hydro2.cpp:1784--1825}

The implementation uses, at the $x$-face between cells $(i,j)$ and
$(i+1,j)$:
\begin{align}
\eta_{\mathrm{face}}
   &= \tfrac12\big[\eta_{i,j} + \eta_{i+1,j}\big], \label{eq:disc_etaf}\\
\dxn\eta\big|_{\mathrm{face}}
   &= (\eta_{i+1,j} - \eta_{i,j}) / \Delta x, \label{eq:disc_dn}\\
\partial_t\eta\big|_{\mathrm{face}}
   &= \tfrac14\big[\eta_{i,j+1}+\eta_{i+1,j+1}-\eta_{i,j-1}-\eta_{i+1,j-1}\big]/\Delta y,
   \label{eq:disc_dt}\\
\grad^2\eta\big|_{\mathrm{face}}
   &= \tfrac12\big[L_9(\eta)_{i,j} + L_9(\eta)_{i+1,j}\big], \label{eq:disc_lap}
\end{align}
where $L_9$ is the 9-point isotropic Laplacian
$L_9(\eta)_{i,j} = \big[4(\eta_{i\pm1,j}+\eta_{i,j\pm1}) + (\eta_{i\pm1,j\pm1}) - 20\eta_{i,j}\big] / (6\Delta x^2)$
(\texttt{Hydro2.cpp:1738--1744}).

For the IC tanh $\eta(x)=\tfrac12[1+\tanh(x/\epsilon)]$, each operator
has the following continuum-vs-discrete behavior on the band:
\begin{itemize}
\item \eqref{eq:disc_etaf}: $\eta_{\mathrm{face}} = \eta(x_{\mathrm{face}}) + \tfrac18\Delta x^2\,\eta''(x_{\mathrm{face}}) + O(\Delta x^4)$. The leading correction is $O(\Delta x^2/\epsilon^2)$.
\item \eqref{eq:disc_dn}: $(\eta_{i+1}-\eta_i)/\Delta x = \eta'(x_{\mathrm{face}}) + \tfrac{1}{24}\Delta x^2\,\eta'''(x_{\mathrm{face}}) + O(\Delta x^4)$, also $O(\Delta x^2/\epsilon^2)$.
\item \eqref{eq:disc_lap}: $L_9(\eta) = \eta''(x) + O(\Delta x^4)$ \emph{when} $\eta''$ is along the grid axis (1D test). 4th-order accurate.
\end{itemize}

The mismatch in truncation order between \eqref{eq:disc_etaf}--\eqref{eq:disc_dn}
($O(\Delta x^2)$) and \eqref{eq:disc_lap} ($O(\Delta x^4)$) is the key
ingredient of (D1): the equipartition relation \eqref{eq:equipart}, which
the continuum analysis uses to cancel
$-\tfrac12\kappa|\grad\eta|^2 + \lambda W = 0$, does not hold discretely
to the same order on both sides. The face evaluation of $|\grad\eta|^2$
carries $O(\Delta x^2/\epsilon^2)$ residue; the evaluation of
$\lambda W(\eta_{\mathrm{face}})$ carries a different $O(\Delta x^2/\epsilon^2)$
residue from the face-averaging of $\eta$.

A small calculation: for $\eta = \tfrac12[1+\tanh(\xi)]$ with $\xi=x/\epsilon$,
\begin{align}
\eta'(\xi) &= \tfrac12\,\mathrm{sech}^2\xi = 2\eta(1-\eta), \\
\eta''(\xi) &= -\mathrm{sech}^2\xi\,\tanh\xi = -2(2\eta-1)\eta(1-\eta), \\
\tfrac12\kappa(\eta')^2 - \lambda W(\eta)
   &= 2\kappa\eta^2(1-\eta)^2/\epsilon^2 - \lambda\eta^2(1-\eta)^2
   \;\stackrel{\eqref{eq:eq_relation}}{=}\; 0,
\end{align}
so equipartition holds in continuum. Discretely, replacing $\eta'$ by
$(\eta_{i+1}-\eta_i)/\Delta x$ introduces a relative error of order
$\Delta x^2/\epsilon^2$ on $(\eta')^2$, while the face-averaged
$\eta_{\mathrm{face}}$ enters $W(\eta_{\mathrm{face}})$ with a different
$\Delta x^2/\epsilon^2$ error.

For the FLAMES \texttt{input\_acoustic\_1d} parameters
($\Delta x \approx 1.95\!\times\!10^{-3}$, $\epsilon=10^{-2}$,
$\Delta x/\epsilon \approx 0.195$), the per-face discrete $\SK^{(B,\mathrm{wb})}_{nn}$ residual scales as
\begin{equation}
|\SK^{\mathrm{disc}}_{nn}|
 \sim \kappa\,(\eta')^2\,(\Delta x/\epsilon)^2
 \sim \kappa\,\cdot\,(1/2\epsilon)^2\,\cdot\,(\Delta x/\epsilon)^2
 = \frac{\kappa\,\Delta x^2}{4\epsilon^4}.
\label{eq:dscale}
\end{equation}
Numerically: $\kappa=6.25\!\times\!10^{-4}$, $\Delta x^2 = 3.8\!\times\!10^{-6}$,
$\epsilon^4 = 10^{-8}$, giving $|\SK^{\mathrm{disc}}_{nn}|$ of order
$0.06$ Pa per face. This is too small by itself to explain a $\sim 30$ Pa
pressure dip. \textbf{However}, the energy update couples to the wave
velocity $S_*$ at line~283 of \texttt{HLLC.H}; the
\emph{state-dependent energy injection} from a $0.06$ Pa
$\SK^{\mathrm{disc}}_{nn}$ residue times an $O(1)$ contact velocity (in
the dimensionless interpretation of the pulse) and integrated over the
$\sim 10^3$-step interaction window is plausibly the right order of
magnitude.

\subsection{Comparison with Form~A as a regression target}

If we were to implement Form~A instead (which is what our \S5 pillbox
argument actually uses), the face stress would be:
\begin{equation}
\SK^{(A,\mathrm{wb})}_{nn} = -\kappa(\dxn\eta)^2 + \tfrac12\kappa|\grad\eta|^2 + \lambda W(\eta_{\mathrm{face}}).
\label{eq:Afaces}
\end{equation}
Compared to Form~B's face stress
\begin{equation}
\SK^{(B,\mathrm{wb})}_{nn} = -\kappa(\dxn\eta)^2 + \kappa\,\eta_{\mathrm{face}}\,L_9(\eta) + \tfrac12\kappa|\grad\eta|^2 - \lambda[\eta W' - W](\eta_{\mathrm{face}}),
\label{eq:Bfaces}
\end{equation}
Form~A:
\begin{itemize}
\item Does \emph{not} use $L_9(\eta)$ at the face (no $\eta\grad^2\eta$ term). The face stencil drops from a 9-point isotropic Laplacian footprint (i.e.\ a $4\!\times\!4$ stencil for the face plus its $\pm1$ neighbors) to a 2-point face-centered difference plus a face-averaged $\eta$. This is a strictly smaller footprint.
\item Discrete equipartition between $(\dxn\eta)^2$ and $W(\eta_{\mathrm{face}})$ requires only that two
$O(\Delta x^2/\epsilon^2)$-truncated quantities combine correctly. There is one less floating-point cancellation than Form~B.
\item Has off-equilibrium behavior differing from Form~B by $-\eta\mu_{\mathrm{chem}}\,\II$ (Eq.\ \ref{eq:diff}). For a pure planar Riemann test ($\mu_{\mathrm{chem}}=0$ at the IC and remains 0 because $\eta$ doesn't advect to a new equilibrium), this difference is purely a small $O(\epsilon)$ correction; for phase-change runs where $\mu_{\mathrm{chem}}\ne 0$ off equilibrium, the two forms genuinely encode different physics.
\end{itemize}

\section{Conclusions and recommended next steps}

\paragraph{The two forms.} Form~A
$\SK^{(A,\mathrm{wb})} = \kappa[\tfrac12|\grad\eta|^2\II - \grad\eta\otimes\grad\eta] + \lambda W(\eta)\II$
and Form~B
$\SK^{(B,\mathrm{wb})} = \kappa[\tfrac12|\grad\eta|^2\II - \grad\eta\otimes\grad\eta] + [\kappa\eta\grad^2\eta - \lambda(\eta W'-W)]\II$
are both well-balanced at planar CH equilibrium. They differ
off-equilibrium by $-\eta\mu_{\mathrm{chem}}\II$ \eqref{eq:diff}. In 1D
both give pointwise $\SK_{nn}=0$ on the planar tanh via equipartition
\eqref{eq:An0}--\eqref{eq:Bn0}.

\paragraph{What the implementation does.} \texttt{HLLC.H:276} together
with \texttt{Hydro2.cpp:1780--1803} implements Form~B with the trace
correction $\Wb^{(B)}(\eta)=\lambda[\eta W' - W]$. The sign is correct.
The construction is consistent. There is \emph{no sign error}.

\paragraph{The likely source of the dip.}
Discrete equipartition fails at $O(\Delta x^2/\epsilon^2)$ because the
9-point Laplacian (4th order), the two-point first-difference (2nd
order), and the face-averaged $\eta$ (2nd order) have differently
ordered leading errors. The Form~B face stress involves all three
operators \eqref{eq:Bfaces}, so its discrete residual on the IC tanh is
the sum of three distinct $O(\Delta x^2/\epsilon^2)$ contributions
that do not cancel by construction. Form~A \eqref{eq:Afaces} would
involve only two of those three operators and has fewer floating-point
cancellations to satisfy.

\paragraph{Proposed numerical diagnostics.}
\begin{enumerate}
\item[(N1)] At $t=0$, dump the per-face $\SK^{(B,\mathrm{wb})}_{nn}$ values for every face in the band region of \texttt{input\_acoustic\_1d}. Plot them. If they are not zero to machine precision, (D1)/(D2) is confirmed.
\item[(N2)] Repeat (N1) under Form~A (compile-time switch or runtime flag). Compare the residuals.
\item[(N3)] Compute the implied cell-divergence
$(\SK_{nn,i+1/2} - \SK_{nn,i-1/2})/\Delta x$ at every cell on the band. The integral of this over the band should be zero (well-balanced), but individual cells will have non-zero residuals.
\item[(N4)] If (N1)--(N3) show non-trivial residuals, the dip-eliminating fix is one of: (i) use Form~A, (ii) tighten the discrete equipartition by matching truncation orders (use a 2nd-order Laplacian to match the 2nd-order first-difference, or a higher-order first-difference), or (iii) post-correct $L_9(\eta)$ on the IC at $t=0$ so the discrete equipartition is exact.
\end{enumerate}

\paragraph{Recommended order of work.}
\begin{enumerate}
\item Run (N1) and (N3) on \texttt{input\_acoustic\_1d} to quantify the discrete residual under Form~B.
\item Implement Form~A as an alternate path (flag-selected) and rerun (N1)/(N3) to compare residuals.
\item Whichever form gives smaller residuals AND continues to give the correct sharp-limit pressure jump in a static-droplet curvature test should be the production form.
\item Document the choice in this writeup and in the \texttt{HLLC.H} comment block at line~250.
\end{enumerate}

\end{document}
